\begin{definition}[Residue-field API]\label{mod:localfield-residuefield}
Let $F$ be a nonarchimedean local field and let $\mathcal O_F$ be its
canonical valuation ring.  Define
\[
 k_F=\mathcal O_F/\mathfrak p_F
\]
using Mathlib's local-ring residue field, and let
$\operatorname{res}_F:\mathcal O_F\to k_F$ be the canonical quotient map.
This map is surjective, and for $x,y\in\mathcal O_F$,
\[
 \operatorname{res}_F(x)=0\iff x\in\mathfrak p_F^1,
 \qquad
 \operatorname{res}_F(x)=\operatorname{res}_F(y)
   \iff x-y\in\mathfrak p_F^1.
\]
The declaration \texttt{reduce} applies this map to an element of $F$
together with a proof that it lies in $\mathfrak p_F^0$; two such integral
elements have the same reduction if and only if they are congruent at depth
one.  Thus the result is independent of the supplied integrality proofs.

Mathlib supplies a \texttt{Finite} instance for $k_F$.
The declaration \texttt{residueFieldFintype} chooses a
noncomputable finite enumeration, and
\[
 q_F=\texttt{residueCard}(F)=|k_F|>1.
\]
Reduction on units is the canonical homomorphism
\[
 \mathcal O_F^\times\longrightarrow k_F^\times.
\]
It is surjective, its underlying residue class is obtained by
$\operatorname{res}_F$, and its kernel is exactly $U_F^1$, represented as
\texttt{unitFiltrationInside} inside $U_F^0$.

The declaration \texttt{teichmuller} specializes Mathlib's adic
Teichmuller map to $\mathcal O_F$.  It is a monoid-with-zero homomorphism
\[
 [\,\cdot\,]_F:k_F\longrightarrow\mathcal O_F
\]
which is a section of reduction.  Consequently it is injective, it sends an
element to zero exactly when that element is zero, and every value satisfies
\[
 [x]_F^{q_F}=[x]_F.
\]
The restriction to $k_F^\times$ is also provided as a homomorphism into
$\mathcal O_F^\times$.

For a local homomorphism $A\to B$ of commutative local rings,
\texttt{residueMapOfLocalHom} is the induced residue-field homomorphism.  It
commutes with the two quotient maps, sends the identity to the identity, and
respects composition.

Finally, for a finite field extension $K/k$, define
\[
 \texttt{residueTrace}_{K/k}=\operatorname{Tr}_{K/k}:K\to k,
 \qquad
 \texttt{residueNorm}_{K/k}=N_{K/k}:K\to k
\]
using Mathlib's linear trace and multiplicative norm.  If $d=[K:k]$, then
for $a\in k$,
\[
 \operatorname{Tr}_{K/k}(a)=d\,a,
 \qquad N_{K/k}(a)=a^d.
\]
For $x\in K$ their images under $k\hookrightarrow K$ are the exact
Frobenius formulas
\[
 \operatorname{Tr}_{K/k}(x)=\sum_{i=0}^{d-1}x^{|k|^i},
 \qquad
 N_{K/k}(x)=\prod_{i=0}^{d-1}x^{|k|^i},
\]
where the displayed equalities are read in $K$.  Both trace and norm compose
in every tower of finite fields.
\LeanFile{LanglandsFirstMainLemma/LocalField/ResidueField.lean}
\lean{LanglandsFirstMainLemma.ResidueField}
\lean{LanglandsFirstMainLemma.residueMap}
\lean{LanglandsFirstMainLemma.residueMap_eq_zero_iff}
\lean{LanglandsFirstMainLemma.residueMap_eq_residueMap_iff}
\lean{LanglandsFirstMainLemma.reduce}
\lean{LanglandsFirstMainLemma.reduce_eq_reduce_iff}
\lean{LanglandsFirstMainLemma.residueFieldFintype}
\lean{LanglandsFirstMainLemma.residueCard}
\lean{LanglandsFirstMainLemma.residueUnits}
\lean{LanglandsFirstMainLemma.residueUnits_ker}
\lean{LanglandsFirstMainLemma.teichmuller}
\lean{LanglandsFirstMainLemma.residueMap_teichmuller}
\lean{LanglandsFirstMainLemma.teichmuller_pow_residueCard}
\lean{LanglandsFirstMainLemma.teichmullerUnits}
\lean{LanglandsFirstMainLemma.residueMapOfLocalHom}
\lean{LanglandsFirstMainLemma.residueTrace}
\lean{LanglandsFirstMainLemma.residueNorm}
\lean{LanglandsFirstMainLemma.residueTrace_trans}
\lean{LanglandsFirstMainLemma.residueNorm_trans}
\leanok
\SourceText{Local notation; Sections Finite-field identities and The non-stable wild calculation.}
\uses{mod:localfield-unitfiltration}
\FormalizationContract
\end{definition}
